% !TEX TS-program = lualatex
В данной главе были рассмотрены популярные подходы к созданию искусственного интеллекта для игр, в частности, игры <<Нарды>>. Традиционные методы, основанные на жёстко заданных правилах и полном переборе возможных действий, просты в реализации, однако не гарантируют принятия решений, стабильно улучшающих положение игрока на каждом ходу.

С другой стороны, подходы, использующие деревья решений, обеспечивающие более высокий уровень интеллекта, требуют значительных вычислительных ресурсов, что ограничивает их практическое применение на современных компьютерах. В этом контексте метод временных разниц представляет собой жизнеспособный компромисс между точностью, скоростью и сложностью, предлагая эффективный баланс для разработки игрового искусственного интеллекта. В связи с этим в данной работе рассматривается техническая реализация алгоритма TD-Gammon.
